\section{Related Work}
\subsection{Authorization-aware RAG}
Chen \emph{et al.} \cite{chen2025sac} integrate access control into a patient-profile RAG prototype and evaluate 40 generated profiles. Their proof of concept demonstrates feasibility but does not separate retrieval exposure from generated disclosure or study the effect of authorization location at the scale used here. Jeong and Lee \cite{jeong2025permission} validate permissions against native IAM endpoints in multi-resource environments and quantify policy-enforcement and answer-quality effects. Chen and Taipalus \cite{chen2026metadata} use metadata both for filtering and fine-grained RBAC, emphasizing secure and relevant retrieval. Together, these studies establish that permission-aware retrieval is an active design direction; our novelty is therefore not the existence of pre-retrieval filtering itself, but a security-ablation methodology that isolates enforcement location, distinguishes context exposure from output disclosure, resolves a retrieval-depth confound, and examines fixed versus proportional authorization scopes at larger corpus sizes.

\subsection{RAG datastore extraction}
Qi \emph{et al.} show scalable extraction from retrieval-in-context RAG systems using instruction-following behavior across multiple model families \cite{qi2025spill}. Di Maio \emph{et al.} propose an adaptive black-box attack that iteratively discovers anchors correlated with previously extracted knowledge \cite{dimaio2025pirates}. These attacks reinforce why retrieved context itself is security-relevant. Our benchmark is not an attempt to reproduce these full adaptive attacks; instead, it uses controlled unauthorized and adversarial prompt families to test whether the retrieval boundary can cross an ACL and whether a narrow marker can reach generated output.

\subsection{Synthetic healthcare records and formalization}
We use Synthea because it produces distributable synthetic health records without introducing real-patient privacy risk \cite{walonoski2018synthea}. Formal analysis uses NuSMV \cite{cimatti2002nusmv} over an abstract state machine. Fuzzy control follows the general notion of graded membership introduced by Zadeh \cite{zadeh1965fuzzy}, but we treat it as one policy implementation rather than as a security guarantee in itself.
